\documentclass[11pt]{article}
\usepackage{amsmath}
\usepackage{geometry}
\geometry{a4paper, margin=1in}

\title{IoT Challenge 1 Ex-4}
\author{Raul Agolli}
\date{March 2026}

\begin{document}

\maketitle

\section*{Q1. Communication Energy Computation}

\begin{itemize}
    \item Packet size ($k$): 2000 bits
    \item Base energy ($E_c$): 300 nJ/bit
    \item Amplifier energy ($\varepsilon$): 1 nJ/bit/m$^2$
    \item Total distance ($D$): 40 m
\end{itemize}

Where:
\begin{align*}
    \text{Transmission: } & E_{tx}(k,d) = kE_c + k\varepsilon d^2 \\
    \text{Reception: } & E_{rx}(k) = kE_c
\end{align*}

\subsection*{(a) Direct Transmission}
In a direct transmission, the sensor sends the message straight to the sink over the full distance of \textbf{40 m}. 
\begin{itemize}
    \item \textbf{1 Transmission} (Sensor $\rightarrow$ Sink) over 40 m
    \item \textbf{1 Reception} (at the Sink)
\end{itemize}

\textbf{Transmission Energy:}
\begin{align*}
    E_{tx} &= (2000 \cdot 300) + (2000 \cdot 1 \cdot 40^2) \\
    E_{tx} &= 600,000 + (2000 \cdot 1600) \\
    E_{tx} &= 600,000 + 3,200,000 \\
    E_{tx} &= 3,800,000 \text{ nJ} = \mathbf{3.8 \text{ mJ}}
\end{align*}

\textbf{Reception Energy:}
\begin{align*}
    E_{rx} &= 2000 \cdot 300 = 600,000 \text{ nJ} = \mathbf{0.6 \text{ mJ}}
\end{align*}

\textbf{Total Energy (Direct):}
\begin{align*}
    E_{total} &= 3.8 \text{ mJ} + 0.6 \text{ mJ} = \mathbf{4.4 \text{ mJ}}
\end{align*}

\subsection*{(b) Multi-hop Transmission}
Based on the figure, there are 4 equal segments (Sensor $\rightarrow$ B1 $\rightarrow$ B2 $\rightarrow$ B3 $\rightarrow$ Sink). This means there are \textbf{4 hops}, each covering a distance of \textbf{10 m} ($40 \text{ m} / 4$).
\begin{itemize}
    \item \textbf{4 Transmissions} (over 10 m each)
    \item \textbf{4 Receptions}
\end{itemize}

\textbf{Transmission Energy per hop:}
\begin{align*}
    E_{tx} &= (2000 \cdot 300) + (2000 \cdot 1 \cdot 10^2) \\
    E_{tx} &= 600,000 + (2000 \cdot 100) \\
    E_{tx} &= 600,000 + 200,000 \\
    E_{tx} &= 800,000 \text{ nJ} = \mathbf{0.8 \text{ mJ}}
\end{align*}

\textbf{Reception Energy per hop:}
\begin{align*}
    E_{rx} &= 2000 \cdot 300 = 600,000 \text{ nJ} = \mathbf{0.6 \text{ mJ}}
\end{align*}

\textbf{Energy per hop:}
\begin{align*}
    E_{hop} &= 0.8 \text{ mJ} + 0.6 \text{ mJ} = \mathbf{1.4 \text{ mJ}}
\end{align*}

\textbf{Total Energy (Multi-hop):}
\begin{align*}
    E_{total} &= 4 \text{ hops} \cdot 1.4 \text{ mJ/hop} = \mathbf{5.6 \text{ mJ}}
\end{align*}

\vspace{0.5cm}
\hrule
\vspace{0.5cm}

\section*{Q2. Strategy \& Justification}

\begin{itemize}
    \item \textbf{Which strategy would you choose for long-term operation?}\\
    \textbf{Multi-hop transmission.}

    \item \textbf{Which performance metric are you optimizing?}\\
    \textbf{Network Lifetime} (specifically, the time until the first node dies, which depends on the lifetime of the battery-powered nodes).

    \item \textbf{Brief Justification:}\\
    While direct transmission uses less \textit{total} network energy (\textbf{4.4 mJ} vs \textbf{5.6 mJ}), it places a huge burden solely on the Motion Sensor (\textbf{3.8 mJ} per transmission). 
    
    The manufacturer specifications indicate the hardware differences: an Aqara Motion Sensor is a battery-powered device, whereas the Philips Hue smart bulbs (B1, B2, B3) and the Sink are plugged into the grid. By choosing the multi-hop strategy, the battery-powered sensor only expends \textbf{0.8 mJ} per transmission, effectively offloading the bulk of the energy cost to the grid-powered bulbs. This significantly prolongs the battery life of the sensor, thereby maximizing the practical lifetime of the network.
\end{itemize}

\end{document}
